\section{Introduction}\label{sec:introduction}

An \emph{oriented graph} is a finite directed graph with no loops, no multiple
arcs, and no pair of opposite arcs.  Its underlying undirected graph need
not be complete or connected.  For a vertex \(v\), write
\[
 N_1^+(v)=\{u:v\to u\},\qquad
 N_2^+(v)=\{w\notin\{v\}\cup N_1^+(v):
                 \text{ there is a vertex }u\text{ with }v\to u\to w\}.
\]
Thus \(N_2^+(v)\) consists of the vertices at directed distance exactly two
from \(v\).  Each vertex is counted once, regardless of the number of paths
reaching it.  Seymour's second neighborhood conjecture asserts that some
vertex satisfies \(\abs{N_1^+(v)}\leq\abs{N_2^+(v)}\).

\begin{theorem}\label{thm:main}
Every nonempty finite oriented graph has a vertex \(v\) such that
\[
 \abs{N_1^+(v)}\leq \abs{N_2^+(v)}.
\]
\end{theorem}

Theorem~\ref{thm:main} resolves the conjecture positively, without a
connectivity or degree hypothesis.  In particular, a sink already satisfies
the conclusion.  The conjecture is a basic comparison between the first two
out-neighborhoods: removing vertices already reached in one step is essential
to its assertion.

The theorem also gives the established vertex- and arc-weighted forms of
the conjecture through Seacrest's equivalences~\cite{Seacrest2015}.
The vertex-weighted conclusions have two different quantifier orders:
for every nonnegative weighting there is a suitable vertex, and there is
a weighting of total mass one for which every vertex is suitable.
Section~\ref{sec:weighted-consequences} states these forms precisely.
Section~\ref{sec:directed-cycle-consequences} records consequences for
short directed cycles.  In particular, a graph with minimum in- and
outdegree at least one third of its order contains a directed triangle.

The conjecture was recorded by Dean and Latka~\cite{DeanLatka1995}.
A \emph{tournament} is an oriented graph with an arc between every pair of
distinct vertices.  This case, also known as Dean's conjecture, was proved
by Fisher~\cite{Fisher1996} using a probability distribution obtained from
Farkas' lemma.  Havet and Thomass\'e gave a combinatorial proof using
local median orders~\cite{HavetThomasse2000}.  Fidler and
Yuster~\cite{FidlerYuster2007} extended the tournament theorem to
tournaments missing a matching, a star, or a clique.  Ghazal corrected the
star-case proof and treated generalized stars~\cite{Ghazal2012}.
Subsequent advances include a matching together with a
star~\cite{DaraFrancisJacobNarayanan2022} and two
stars~\cite{DaamouchAlMninyGhazal2025}.

Other approaches impose degree conditions or weaken the desired ratio.
Kaneko and Locke established the conjecture for minimum outdegree at most
six~\cite{KanekoLocke2001}; a recent computer-assisted preprint treats
minimum outdegree seven~\cite{SadhukhanSandeepSen2026}.
Chen, Shen and Yuster proved a universal second-to-first neighborhood
ratio of approximately $0.657298$~\cite{ChenShenYuster2003}.
Huang and Peng raise this to approximately $0.715538$ and additionally
report a computational bound of $0.74530$~\cite{HuangPeng2024}.
Botler, Moura and Naia~\cite{BotlerMouraNaia2023}, followed by
Espuny D\'iaz, Gir\~ao, Granet and Kronenberg~\cite{EspunyDiazEtAl2025},
studied every orientation of binomial random graphs; the latter result
covers every fixed edge probability below $1/2$, with probability tending
to one as the order tends to infinity.  Brukhman's recent dense-case
preprint concerns graphs of order $n\leq 2\delta+2$, where $\delta$ is
the minimum outdegree~\cite{Brukhman2026}.

Our proof builds on the minimal-counterexample and graph-product
methods of Brantner, Brockman, Kay and
Snively~\cite{BrantnerBrockmanKaySnively2009} and on Seacrest's
subset perspective~\cite{Seacrest2019}.
We adapt the edge-deletion idea in the proof of Seacrest's Lemma~4 to
obtain the stronger subset inequality needed here, and give that argument
in full.  The exact set-image convention and inequality are stated below.
The main additional ingredient is a pruning bound for arbitrary finite
binary relations.  Its deletion cost depends on the points left uncovered
after pruning; this dependence makes the final extremal argument possible.

\paragraph{Proof strategy.}
For a digraph \(D\) and a set \(S\) of its vertices, define
\[
 F_D(S)=\{y:\text{ there is }x\in S\text{ with }x\to y\},
 \qquad F_D^2(S)=F_D(F_D(S)).
\]
The operator \(F_D\) does not subtract \(S\).  In an oriented graph there is
no directed walk of length two from a vertex to itself, and hence
\(N_2^+(v)=F_D^2(\{v\})\setminus F_D(\{v\})\).

We first suppose that a counterexample exists and choose one with minimum
order and then minimum number of arcs.  Deleting carefully chosen arcs
shows that every nonempty proper vertex set \(S\) satisfies
\[
 \abs{F_D^2(S)\setminus F_D(S)}
 <
 \abs{F_D(S)\setminus S}.
\]
Replacing each vertex by a sufficiently large transitive tournament turns
this into a strict inequality
\begin{equation}\label{eq:overview-concavity}
 \abs{U}+\abs{F_G^2(U)}<2\abs{F_G(U)}
\end{equation}
for every nonempty proper subset \(U\) of the new oriented graph \(G\).
Every vertex of \(G\) has positive indegree.  These reductions are proved in
Section~\ref{sec:reduction}.

The obstruction to~\eqref{eq:overview-concavity} is an extremal argument
on two families of ordered pairs in $V(G)\times V(G)$.  A coordinate image
advances that coordinate along an arc of $G$ and leaves the other coordinate
fixed.  The families are compatible when the first-coordinate image of the left family and the
second-coordinate image of the right family are disjoint.  A point in
neither image is called uncovered.  Require both families to contain the
diagonal, write $M$ for their combined size and $d$ for the number of
uncovered points, and maximize $M+d$ over all such compatible pairs.
The columns of the left family and the rows of the right family are
nonempty proper vertex sets, so
\eqref{eq:overview-concavity} applies to each of them.  Taking complements
of the second iterated images then produces prospective families with
combined size greater than $M+2d$.

These families may conflict, but their diagonal members are still present
and conflict-free.  The pruning lemma in
Section~\ref{sec:pruning}, valid for every finite binary relation,
restores compatibility and preserves their diagonal members.
In this application its two deletion-cost terms are at most $d$ and $d'$,
where $d'$ is the number of points left uncovered by the retained families.
Thus at most $d+d'$ members are deleted, and the new extremal objective is
strictly greater than
\[
 M+2d-(d+d')+d'=M+d.
\]
Section~\ref{sec:incompatibility} proves these bounds and obtains the
contradiction.  This is why the pruning lemma must charge part of its
cost to the points left uncovered after deletion.

The pruning proof bounds a matching in an augmented bipartite graph.
Changing the two coordinates of a pair in opposite orders factors the
same matrix through two different index sets attached to a conflict.
A simultaneous maximum-rank choice of the entries makes the kernels of
these maps control the matching size.  The full argument is independent
of the orientation assumption and is given in Section~\ref{sec:pruning}.
